% Author's solution, images/HW04/IMG_9189.HEIC--IMG_9193.HEIC.
% IMG_9192 replaces IMG_9191's preliminary e_2,e_3 choices.
% The assignment and the completed equations use L:R^3 -> R^2 and x+y+z.
% Describe the supplied solution of the full kernel system, without
% retaining the invalid IMG_9190 argument restricted to scalar t=0.
\noindent\textbf{Solution.}
\begin{enumerate}[(a)][2]
  \item The condition \(L(x,y,z)=0\) gives the system
    \[
      x+y+z=0,
      \qquad
      x+2y+3z=0.
    \]
    Subtracting the first equation from the second gives \(y+2z=0\).
    Hence \(y=-2z\), and the first equation then gives \(x=z\).
    The solutions are exactly
    \[
      (x,y,z)=(z,-2z,z)=z(1,-2,1).
    \]
    Therefore
    \[
      \ker L=\operatorname{span}\{(1,-2,1)\}.
    \]
    The equation \(a(1,-2,1)=0\) forces \(a=0\), so this single
    vector is linearly independent and forms a basis of the kernel.

  \item Choose the ordered list
    \[
      E=(e_1,e_2,e_3)
        =\bigl((1,-2,1),(1,0,0),(0,1,0)\bigr).
    \]
    If \(ae_1+be_2+ce_3=0\), then
    \[
      (a+b,-2a+c,a)=(0,0,0),
    \]
    giving \(a=b=c=0\). Thus these vectors are linearly independent
    and form a basis of \(\mathbb{R}^3\), with the kernel basis first.

  \item Set
    \[
      f_1=L(e_2)=(1,1),
      \qquad
      f_2=L(e_3)=(1,2).
    \]
    If \(af_1+bf_2=0\), then
    \[
      a+b=0,
      \qquad
      a+2b=0.
    \]
    These equations imply \(a=b=0\). Hence
    \(F=(f_1,f_2)=\bigl((1,1),(1,2)\bigr)\) is an ordered basis
    of \(\mathbb{R}^2\).

  \item Since
    \[
      L(e_1)=0,
      \qquad
      L(e_2)=f_1,
      \qquad
      L(e_3)=f_2,
    \]
    their \(F\)-coordinate columns give
    \[
      M=\begin{pmatrix}0&1&0\\0&0&1\end{pmatrix}.
    \]
    \marginnote{Each column of \(M\) records \(L(e_j)\) in the basis
      \(F\). The kernel vector gives the zero first column.}

  \item Applying \(L\) to each column of \(E\) gives
    \[
      L(E)=[\,0\ \ f_1\ \ f_2\,]
        =\begin{pmatrix}0&1&1\\0&1&2\end{pmatrix}.
    \]
    On the other hand,
    \begin{align*}
      FM
        &=[\,f_1\ \ f_2\,]
          \begin{pmatrix}0&1&0\\0&0&1\end{pmatrix}\\
        &=[\,0\ \ f_1\ \ f_2\,]
          =\begin{pmatrix}0&1&1\\0&1&2\end{pmatrix}.
    \end{align*}
    Thus \(L(E)=FM\).
\end{enumerate}
